\chapter{Discussion and Synthesis}
\label{ch:discussion}

The studies of this thesis were presented in sequence, each answering a research
question. This chapter steps back to consider them as a whole, and the whole is an
argument about the adoption of formal methods. It argues that the studies together
remove the principal barrier to that adoption and open a capability that did not
previously exist (Section~\ref{sec:syn-adoption}); it draws out the methodological
thread that runs through all of them, the disciplined separation of heuristic
inference from formal validation (Section~\ref{sec:syn-heuristic}); it examines
what the studies jointly establish about the value of inferred specifications
(Section~\ref{sec:syn-value}); it confronts the tensions and limitations that
cut across the chapters (Section~\ref{sec:syn-tensions}); and it considers the
combined threats to validity that no single chapter's threats section captures
(Section~\ref{sec:syn-threats}).

\section{The Adoption Argument}
\label{sec:syn-adoption}

The thesis began from a puzzle: formal methods are mature and effective, yet used
only in a few niche domains. Its answer is that the binding constraint is the
manual cost of the specifications verification requires, and the four studies form
a single argument that this constraint is substantially removable --- and that
removing it opens a capability beyond cheaper verification.

The argument has two movements. The first establishes that inference removes the
cost: the engine produces accurate specifications (Chapter~\ref{ch:inference}, the
fidelity result), and those specifications carry measurable value to a demanding
consumer (Chapter~\ref{ch:testgen}, the test-generation result). Together these
answer the obvious objection to automatic inference --- that a cheap specification
is a worthless one --- by showing the inferred specifications are both accurate and
useful. If the argument stopped here it would establish a worthwhile but
incremental result: verification made cheaper, which lowers the entry cost to the
niche without necessarily widening it.

The second movement is what bears on adoption most directly, and it is the
thesis's distinctive claim. Specifications inferred and carried at the
\emph{library boundary} (Chapter~\ref{ch:distribution}) make a dependency's
contract available to a consumer who lacks its source --- and because a caller's
specification is \emph{derived from} its callees' (Chapter~\ref{ch:compositional}),
a change to a library propagates into the contracts of the code that uses it. This
is the basis for behavioural-compatibility analysis across library versions: a
capability that manual specification, confined to whole systems specified in
isolation, never made routine, and one that answers a question --- is this upgrade
safe? --- that every practitioner who manages dependencies faces. The adoption
argument is therefore not merely ``verification is now cheaper'' but ``inferred,
propagating, library-boundary specifications offer something practitioners want and
could not previously have,'' which is the kind of new value that moves a technology
out of a niche rather than merely lowering the cost of entering it.

The studies depend on one another in ways the individual chapters make concrete.
The distribution mechanism is useful only because the inference produces something
worth distributing, and its validation deliberately exercises the full
inference-to-embedding path rather than synthetic specifications, to show it
transports what the engine actually emits. The compositional analysis operates on
the specifications the inference produces, and the compatibility capability it
enables depends on the distribution that puts those specifications at the library
boundary. The test-generation value is what makes the whole worth pursuing: without
a demonstrated downstream benefit, the machinery would serve no purpose. The
chapters can be read independently, but they describe one argument --- that the
specification-cost barrier to formal methods is removable, and that removing it at
the library boundary opens compatibility analysis as a new capability.

A consequence worth drawing out is that the pipeline closes a loop the
literature has left open. The verification community has mature tools to
\emph{check} specifications and the testing community has mature tools to
\emph{consume} them, but both have always required the specifications as input,
and the manual cost of supplying that input is the bottleneck the introduction
identified. The pipeline supplies the input automatically and gets it to the
tools that need it, including across the source-to-binary boundary that
separates a library's author from its users. The contribution of the thesis is
not any single stage --- each has precursors --- but the demonstration that the
stages compose into a working whole whose output carries measurable value.

\section{The Coherence of the Whole}
\label{sec:syn-coherence}

Before examining the methodological thread and the cross-cutting threats, it is
worth making explicit what binds the studies into a single argument rather than a
sequence of separable results. The answer is that they share not only a system ---
though they do, every study using or extending JML-Inferrer --- but a single
underlying question, pressed at each stage of the argument: \emph{can the manual
cost of specifications be removed, and removed in a way that opens new capability
rather than merely cheapening an old one?} Inference removes the cost of producing
specifications; validation ensures the cost reduction does not sacrifice
soundness; distribution puts the specifications at the library boundary where
verification across dependencies happens; and the compositional measurement
establishes the caller-on-callee dependence on which behavioural-compatibility
analysis rests. Each study is the same question pressed against a different
obstacle to the proposition that automatic, library-boundary specifications can
move formal methods beyond the niche.

The coherence shows in how the studies' results depend on one another's framing. The
distribution study's insistence on exercising the real inference pipeline, not just
synthetic specifications, is motivated by the need to show it transports what the
inference study actually produces. The compositional study's downstream probe is
framed as a re-run of the test-generation study's experiment, so that its preliminary
signal is measured on the same metric and is interpretable as a continuation of the
value finding. The validation layer described in the inference chapter is the same
layer whose discharge results the test-generation chapter reports and whose
soundness guarantee the distribution chapter relies on when it claims to transport
verified contracts. Remove any one study and the others lose not their individual
results but part of their motivation: inference without value is unmotivated, value
without reach is academic, reach without completeness is premature, and completeness
without the prior three is a refinement of nothing. The thesis is the argument that
the four together establish a proposition none establishes alone, and the coherence
is the dependency among their framings, not merely their shared use of one system.

\section{Heuristic Inference, Formal Validation}
\label{sec:syn-heuristic}

The deepest methodological commitment of the thesis, recurring in every chapter,
is the disciplined separation of heuristic plausibility from formal soundness.
The inference engine is heuristic: it matches syntactic patterns, and a pattern
can fire where it should not. Rather than claim a soundness the engine cannot
deliver, the thesis treats inference and verification as two independent layers
in which the first proposes and the second disposes. This stance --- infer
aggressively, verify conservatively --- is, in a sense, the methodological
thesis underlying the technical one, and it is worth examining why it works.

The alternative stances each fail in an instructive way. A purely heuristic
engine with no verification would emit unsound clauses with no recourse, and its
output could not be trusted by a verifier or a safety-critical consumer. A purely
sound engine --- one that emits only clauses it can prove by construction ---
would be a verification-condition generator in disguise, and would inherit the
scalability problems that make full WP computation impractical: it would need the
loop invariants that are the hardest thing to infer, and it would emit nothing
for the methods where it could not complete the proof. The two-layer
architecture escapes both failure modes. The heuristic layer is free to be
aggressive, emitting any clause a pattern suggests, because the validation layer
will catch the unsound ones; the validation layer is free to be conservative,
rejecting any clause it cannot prove, because the heuristic layer has already
done the creative work of proposing candidates. The division of labour is
between a creative-but-fallible proposer and a rigorous-but-uncreative checker,
and it is the same division that underlies Houdini's guess-and-prune
architecture~\cite{flanagan2001houdini}, here grounded in code structure rather
than in templates.

The independence of the two oracles --- manual validation against documented
intent, formal discharge against the implementation --- is a refinement of the
same theme that Chapter~\ref{ch:testgen} exploits. The two oracles catch
different errors: manual validation catches clauses that are formally consistent
but miss the documented point, and formal discharge catches clauses that look
right to a human reviewer but conflict with the code's actual arithmetic or
aliasing. Reporting them separately, rather than collapsing them into a single
accuracy figure, is what allows the thesis to claim that its specifications are
accurate under \emph{both} a human and a machine oracle --- a stronger claim than
either alone, and one that the discharge rate exceeding the manual-validation
precision (Chapter~\ref{ch:testgen}) makes vivid: the verifier accepts clauses
the humans called incomplete, because completeness against documentation and
consistency with implementation are genuinely different properties.

\section{What the Studies Jointly Establish About Value}
\label{sec:syn-value}

Each study contributes a piece of evidence about the value of inferred
specifications, and the pieces reinforce one another.

Chapter~\ref{ch:testgen} establishes the most direct value claim: supplying
inferred specifications to a language model improves its test generation by a
large and statistically robust margin --- 272\% more tests and 40.7 percentage
points of mutation score over a signature-only baseline --- and the improvement
is not replicated by generic prompt enrichment, which isolates the specification
as the active ingredient. This is value to a consumer, measured by the consumer's
own metric.

Chapter~\ref{ch:distribution} establishes that this value is \emph{reachable}.
A specification that improves test generation is useless to a consumer who
cannot obtain it, and most consumers --- developers integrating a compiled
dependency, models reading an API at inference time --- cannot obtain a
specification that lives only in the source tree. The 100\% lossless embedding at
sub-eleven-percent overhead establishes that the value demonstrated in
Chapter~\ref{ch:testgen} can be delivered to consumers who lack the source, which
is the precondition for the value mattering at scale.

Chapter~\ref{ch:compositional} establishes that there is \emph{more} value
available than the single-pass inference captures. The 42{,}370 clauses the
second pass adds, dominated by shapes the single pass cannot express, are not
hypothetical: they are produced by a working pass on real libraries, and the
preliminary downstream signal --- a positive though noisy effect on test count
when the refined specifications are supplied to the model --- connects them back
to the value metric of Chapter~\ref{ch:testgen}. The chapter establishes a
ceiling that the earlier chapters were operating below, and a means of raising
toward it.

Read together, the three value claims form an argument with a structure: inferred
specifications are valuable (Chapter~\ref{ch:testgen}), the value is deliverable
to those who need it (Chapter~\ref{ch:distribution}), and there is more value to
be had by inferring more thoroughly (Chapter~\ref{ch:compositional}). No single
chapter makes this argument; it emerges only from the conjunction, which is the
sense in which the thesis is more than the sum of its three studies.

The argument's structure also clarifies what would weaken it. If the value claim
failed --- if specifications did not improve the consumer --- the other two would be
machinery serving no purpose, and the thesis would be a well-engineered solution to a
non-problem. If the value held but the reach claim failed --- if specifications could
not be delivered without the source --- the value would be confined to the teams that
already have the source and need the specifications least, and the contribution would
be academic. If value and reach held but the completeness claim showed the single
pass already captured everything --- if the second pass added nothing --- the thesis
would still stand on its first two claims but would lack the forward-looking element
that the gap and its recovery supply. The three claims are thus not merely additive
but mutually supporting: value justifies the machinery, reach makes the value
matter, and completeness shows the value is not yet exhausted. The thesis's strength
is that all three hold, and its argument is that their conjunction establishes the
proposition --- that automatic specifications can be as good as the practice
needs --- that motivates the whole.

\section{The Economics of the Pipeline}
\label{sec:syn-economics}

The thesis's premise is economic --- that the manual cost of specifications is the
binding constraint --- and the synthesis can make the economic argument the studies
jointly support more explicit. The cost of a manually specified codebase scales
with its size: every method needs a contract, each contract costs human effort, and
the total is the per-contract cost times the method count, which for a large
codebase is prohibitive. The pipeline replaces this with a different cost
structure. Inference costs machine time proportional to code size but negligible
human time; validation costs machine time for the solver and human time only for the
flagged residual; distribution costs machine time and no human time; refinement
costs machine time and no human time. The human cost, which dominated the manual
approach and scaled with method count, is reduced to reviewing the flagged clauses
--- a fraction of the total --- and scales with the failure rate rather than the
method count.

This is the economic transformation the thesis demonstrates is achievable. The
discharge results bound the human review cost: with over ninety percent of clauses
discharged automatically and the flagged residual a few percent, the human reviews a
small fraction of the inferred clauses rather than authoring all of them. The
throughput figures bound the machine cost: inference at milliseconds per file,
embedding at thousands of methods per second, refinement at milliseconds per method,
all of which disappear into a build that already takes longer. The net effect is that
the cost of having specifications shifts from a human cost that scales with size to a
machine cost that scales with size plus a human cost that scales with the much
smaller failure rate --- which is precisely the shift that makes specifications
affordable for codebases where manual specification never was.

The economic argument has a limit the thesis is careful to respect. The inferred
specifications are not the equal of expert hand-written ones in completeness --- the
recall gap is real, and a verification effort requiring strong, complete contracts
would still need human authorship for the methods the engine specifies only
partially. The economic transformation is not from manual to free but from manual to
mostly-automatic-with-review, and its value is greatest where partial-but-abundant
specifications suffice --- test generation, lightweight verification, documentation,
model context --- which is exactly the set of consumers the thesis targets. For the
heaviest verification, the pipeline lowers the entry cost without eliminating the
expert effort the last mile requires, and the honest economic claim is that it
shifts the cost structure favourably for most uses, not that it abolishes the cost
for all.

\section{Tensions and Design Trade-offs}
\label{sec:syn-tensions}

Several tensions cut across the chapters, and acknowledging them is part of an
honest account of the work.

The first is the \emph{precision--recall tension} in inference. The engine
favours precision: it emits a clause only when a pattern matches a structure it
trusts, and it declines to guess where it is unsure, which is why a few methods
receive only empty contracts. A more aggressive engine would emit more clauses
and catch more of the documented intent, at the cost of more false positives for
the validation layer to filter. The thesis sits deliberately on the precise side
of this trade-off, because the downstream consumers --- a verifier, a test
generator --- are better served by a smaller set of trustworthy clauses than by a
larger set of doubtful ones, and because the comparison with language-model
inference (Chapter~\ref{ch:testgen}) shows that the imprecise-but-complete
strategy produces exactly the hallucinated clauses the thesis aims to avoid.

The second is the \emph{heuristic-versus-formal tension} in coverage. The
engine's heuristic front end can emit clauses the formal back end cannot
discharge --- the 7.1\% of clauses that hit the three tractability limits of
Chapter~\ref{ch:testgen}. These clauses are not wrong; they are unverifiable at
the current strictness with the current solver, and the thesis retains them as a
separate class (useful as oracle information to a language model, outside the
verified-sound surface) rather than discarding them. The tension is between the
engine's reach and the verifier's grasp, and it is resolved not by restricting
the engine to what the verifier can prove --- which would discard genuinely
useful clauses --- but by labelling each clause with its verification status and
letting the consumer decide how much to trust the unverified ones.

The third is the \emph{single-pass-versus-compositional tension} in
interprocedural analysis, which Chapter~\ref{ch:compositional} makes explicit.
The single pass is fast, simple, and sound for the unconditional case; the
compositional pass is richer but more complex. The thesis does not argue that the
compositional pass should always run --- it is an opt-in enhancement, and for a
consumer that wants only the cheap unconditional obligations the single pass
suffices --- but it does establish that the choice is a real one with measurable
consequences, and that an inferrer wishing to remain single-pass could capture
the easy half of the gap (branch-guard lifting, which needs only the local
control-flow context) without the harder half (polymorphic disjunction, which
needs the class hierarchy). Naming the trade-off precisely, and quantifying what
each side costs and yields, is itself a contribution.

\section{The Language Model as Producer and Consumer}
\label{sec:syn-llm}

A theme that surfaces in two of the studies, and that the synthesis can draw out
more sharply than either study alone, is the asymmetric role of the language
model in the specification life cycle. The model appears twice: as a candidate
\emph{producer} of specifications (the comparison baseline of
Chapter~\ref{ch:testgen}) and as a \emph{consumer} of them (the central
experiment of Chapter~\ref{ch:testgen} and the downstream probe of
Chapter~\ref{ch:compositional}). The two roles produce opposite verdicts, and
the contrast is instructive.

As a producer, the model is weaker than the heuristic engine on the dimensions
that matter for a downstream verifier: it attains lower precision, because it
hallucinates clauses that are plausible in surface form but false against the
code, and far lower run-to-run consistency, because its sampling is stochastic
where the static analysis is deterministic. The model's productions are not
grounded in the code's structure --- it generates a specification the way it
generates any text, by prediction --- and the lack of grounding shows up as
exactly the failure mode a grounded analysis avoids. As a consumer, by contrast,
the model is excellent: given a grounded specification, it produces markedly
better tests, because the specification supplies the statement of intent the
model cannot reliably manufacture for itself. The model that hallucinates when
asked to produce a contract uses a supplied contract faithfully.

The asymmetry has a clean explanation. Producing a sound specification requires
grounding in the code's actual semantics, which the heuristic engine has (it
reads the AST) and the model lacks (it predicts from text). Consuming a
specification requires fluency in turning intent into tests, which the model has
in abundance and the static engine lacks entirely. The two are complementary in
exactly the way the thesis's architecture exploits: a grounded producer feeds a
fluent consumer, and the combination outperforms either alone. The methodological
caution that the thesis maintains --- use the static engine to produce, use the
model to consume --- is not a prejudice against language models but a
consequence of where each kind of system is strong. A future system might close
the producer gap by grounding a model's specification generation in a verifier's
feedback, in the spirit of Houdini's propose-and-prune loop with a model as the
proposer; the thesis does not build that system, but its results delineate the
niche such a system would have to fill.

\section{Toward Automated Specification Engineering}
\label{sec:syn-vision}

Stepping back from the individual results, the thesis can be read as an argument
for a broader proposition: that specifications, long treated as a manually
authored input to the verification and testing pipeline, can become an
automatically maintained \emph{artefact} of it, computed from code, certified
against it, distributed with it, and refined as analysis improves. The four
stages of the pipeline correspond to the four operations one would want on such
an artefact --- create it, check it, ship it, improve it --- and the thesis
demonstrates each on real code.

The proposition matters because it changes the economics that have kept formal
specifications marginal. If a specification is a manual artefact, every method
that lacks one is a method whose specification will probably never be written,
because the marginal cost falls on a human who has other work. If a specification
is an automatically maintained artefact, every method has one by default, kept
current by re-running the inferrer in continuous integration, certified by
re-running the verifier, and shipped by re-running the embedder. The blank page
that the introduction identified as the barrier is filled not once, heroically,
but continuously, mechanically, as a side effect of the build. The cost that
Matichuk et al.~\cite{matichuk_cost_2015} measured as scaling with specification
volume is paid by a machine whose throughput (Chapter~\ref{ch:inference}'s 88
milliseconds per file, Chapter~\ref{ch:distribution}'s thousands of methods per
second) makes the volume irrelevant.

This is, of course, an aspiration that the thesis only partially realises. The
specifications the engine produces are accurate but incomplete; the verifier
discharges most but not all of them; the distribution is lossless but the
inference behind it is heuristic; the refinement recovers compositional
information but its downstream value is so far only a preliminary signal. The
thesis demonstrates that each stage works on well-structured library code, not
that the pipeline is ready to maintain the specifications of an arbitrary
enterprise system. But the demonstration is enough to establish the proposition's
plausibility: the operations compose, they run at practical speed on real
libraries, and their output carries measurable value. Whether the proposition
scales to the full diversity of production software --- the application code, the
framework code, the concurrent and reflective and dynamically-loaded code that
the controlled corpora exclude --- is the question that the thesis's limitations
mark out as the frontier, and that its future work proposes to approach.

\section{Combined Threats to Validity}
\label{sec:syn-threats}

Each chapter states the threats specific to its study; this section addresses the
threats that span them, which no single chapter's threats section captures.

The most important cross-cutting threat is \emph{corpus homogeneity}. Apache
Commons Lang appears in all three studies, Commons IO in two, and the libraries
are uniformly mature, well-structured, open-source utility and value-type code.
This is a deliberate choice --- using shared corpora keeps cross-chapter
comparisons free of corpus-difference confounds, and mature libraries are the
realistic target for specification inference --- but it bounds external validity
in a way that compounds across chapters. A result that holds on Commons Lang in
all three studies is well-established on Commons Lang and correspondingly less
established elsewhere. Chapter~\ref{ch:compositional} mitigates this most
directly, deliberately spanning four idiom families including generics-heavy and
numerical code and finding the result holds across all five libraries; but
application code, framework code, and proprietary enterprise code remain
unrepresented in every study, and the \emph{other} method category --- event
handlers and callbacks --- is empty in the controlled corpus of
Chapter~\ref{ch:testgen}. Generalisation beyond well-structured library code is
the single largest open question the thesis leaves, and it is named as such in
every chapter that touches it.

The second cross-cutting threat concerns the \emph{language model}. The
test-generation study (Chapter~\ref{ch:testgen}) and the downstream probe
(Chapter~\ref{ch:compositional}) both use a single model family, for budget
reasons, and both defend an \emph{ordering} of conditions rather than absolute
scores. The defence rests on the model-independence of the JML format and on
multi-model studies in the literature reporting consistent qualitative findings
across model families, but a full multi-model replication at the scale of the
thesis's experiments is not performed and is named as future work. The risk is
that the magnitude of the effects is model-specific even if their direction is
not; the thesis claims only the direction.

The third cross-cutting threat is the \emph{heuristic-soundness} concern that the
two-layer architecture is designed to address but cannot fully eliminate. Every
clause is discharged through OpenJML, but discharge is itself bounded by the
solver's tractability, and the unsupported clauses (Chapter~\ref{ch:testgen}) are
neither proved nor refuted --- they are retained on the strength of manual review
and their usefulness as oracle information, outside the verified surface. A
consumer that requires every clause to be machine-verified would discard these,
and the thesis's claim that its specifications are accurate must be read with the
distinction between verified, flagged, and unsupported clauses kept in view. The
architecture reduces the risk that unsound clauses reach a consumer; it does not
reduce it to zero, and the honest statement is that the residual risk is
concentrated in exactly the clauses the chapter labels as unsupported.

The fourth, specific to the most recent work, is that the downstream signal of
Chapter~\ref{ch:compositional} measures test count rather than mutation kills,
because the generated tests did not all compile cleanly. This is the weakest
empirical link in the thesis, and it is reported as a preliminary signal rather
than a result for exactly that reason. The fuller mutation-coverage comparison
between baseline and refined specifications is the natural completion of the
argument and the top-priority follow-up; until it is done, the claim that the
compositional additions carry downstream value rests on a non-zero but noisy
test-count delta rather than on the fault-detection metric the rest of the thesis
relies on.

\section{Engineering Lessons}
\label{sec:syn-engineering}

Beyond its empirical findings, the thesis accumulated engineering lessons that
recur across the studies and that are worth stating as guidance, because they are
the kind of knowledge that a thesis organised around results does not otherwise
record.

The first is that a high-but-imperfect aggregate should be disaggregated before it
is accepted. The interface-method embedding bug of Chapter~\ref{ch:distribution}
hid behind a 96.9\% roundtrip rate that looked like success; disaggregating the
failures by program element revealed that every one was an interface method, which
pointed directly at the missing lifecycle callback. The same discipline applies to
the discharge results of Chapter~\ref{ch:testgen}, where disaggregating the
non-discharged clauses by structural pattern revealed the three tractability limits
rather than an undifferentiated failure rate. The lesson is that an aggregate
metric is a starting point for investigation, not an endpoint, and that the
residual in any imperfect measurement usually has a specific, findable cause.

The second is that the cases a system handles silently are the cases most likely to
harbour bugs, precisely because they generate no signal when they fail. Interface
methods received no embedded annotation and threw no error; the failure was visible
only in the aggregate roundtrip rate, not in any exception or log. A test suite that
asserted the presence of an embedded annotation on every method, including those
without bodies, would have caught the bug immediately. The lesson is to test the
silent cases explicitly, because the loud cases test themselves.

The third is that a deliberately conservative policy --- rejecting unsound
substitutions, gating clauses by return type and termination, declining to infer
where uncertain --- costs recall but buys a guarantee worth more than the recall it
costs. Across all three studies the engine favours the conservative choice, and the
consequence is a lower-bound character to its claims: the inferred specifications
are at least as accurate as reported, the compositional gain is at least as large,
the embedding loses at most nothing. The lesson is that for a tool whose output
feeds a verifier or a test generator, the conservative bias is the right one,
because a missed clause is recoverable by a later, more aggressive analysis while
an unsound clause that reaches a consumer is not.

\section{The Recurring Role of Negative Results}
\label{sec:syn-negative}

A feature of the thesis that the synthesis can highlight is the consistent value
of its negative and surprising results, which arise directly from the commitment
to measurement on real code. The studies are not a sequence of confirmations;
they include findings that ran against the initial expectation, and these are
among the most informative.

The clearest is the per-spec compression result of
Chapter~\ref{ch:distribution}. The intuition that compressing each
specification would reduce the embedding overhead is natural, and the measurement
refutes it decisively: per-spec DEFLATE made the overhead dramatically worse,
because the payloads are too small for the compressor to find redundancy and the
compressed payloads forfeit the constant-pool deduplication the uncompressed
format enjoys. The negative result is informative beyond its immediate context:
it shows that the right granularity for any compression of embedded
specifications is the class file, not the individual specification, because only
at the class-file level is there enough shared structure to amortise. A study that
reported only the successful optimisations would have left a deployer free to try
per-spec compression and rediscover the failure.

A second is the P2-versus-P3 inversion of Chapter~\ref{ch:testgen}: the guidance
condition produces more tests but a lower mutation score than the specification
condition. The result was the point of including the guidance control, but its
magnitude --- that more tests can mean worse fault detection --- is a finding with
methodological reach, cautioning the entire test-generation field against test
count as a quality proxy. A third is the precise characterisation of the three
tractability limits: rather than reporting an undifferentiated discharge-failure
rate, the study identifies the structural patterns that defeat the solver, turning
a negative result (some clauses do not discharge) into positive knowledge (these
specific shapes exceed these specific solver capabilities). A fourth is the
Guava-overhead inversion, where real inferred specifications produced \emph{lower}
overhead than synthetic ones on the largest library, contrary to the expectation
that real specifications are richer and therefore costlier --- explained by the
larger denominator of Guava's dense generic classes.

These results share a provenance: they came from measuring rather than assuming,
on real artefacts whose properties did not match the intuition. The thesis's
commitment to measurement on production code is what surfaced them, and their
presence is evidence that the commitment paid off --- a study that argued from
construction would have asserted the intuitions and missed the corrections.

\section{What Would Falsify the Thesis}

A claim worth defending is a claim that could in principle be refuted, and stating
what would falsify the thesis's central claim sharpens what the evidence does and
does not establish. The thesis claims that heuristic inference, validated against a
formal oracle, produces specifications accurate enough to be useful, distributable,
and valuable to consumers. Several findings would have refuted it, and the studies
were designed so that they could have.

The fidelity claim would have been refuted by a low discharge rate or a low manual-
validation precision --- if the verifier had rejected most inferred clauses, or if
the raters had judged most of them incorrect, the specifications would not be
accurate enough to be useful. The observed rates (above 90\% on both oracles) do not
refute it, but they could have; an engine whose patterns fired indiscriminately
would have produced exactly the low rates that would falsify the claim. The value
claim would have been refuted by the guidance control: if the guidance condition had
matched or exceeded the specification condition on mutation score, the effect would
have been attributable to prompt enrichment in general rather than to specifications
in particular, and the value of specifications specifically would be unestablished.
The control did not refute the claim --- guidance produced more tests but a lower
mutation score --- but it was included precisely because it could have, and a
weaker experimental design that omitted it would have left the claim vulnerable to
exactly that alternative explanation.

The reach claim would have been refuted by a roundtrip rate below 100\% that
resisted explanation, or by an overhead so large as to be impractical; the
interface-method bug produced a sub-100\% rate that, had it been accepted rather
than disaggregated, would have left the lossless claim unproven. The completeness
claim would have been refuted by a second pass that added few clauses, or clauses
of shapes the single pass could already produce --- if the refinement had merely
duplicated the single pass's output, it would have shown the single pass already
compositional. The observed additions (tens of thousands of clauses, dominated by
shapes the single pass cannot express) do not refute it, but a single pass that was
already capturing the guarded and polymorphic obligations would have produced the
small, redundant additions that would falsify the claim.

That each claim had an identifiable way to fail, built into the design of the study
that tested it, is what distinguishes the thesis's empirical assertions from mere
descriptions of a system that works. The claims survive not because they could not
have failed but because they were tested in ways that would have exposed them if
they were false, and the surviving claims are correspondingly load-bearing.

\section{Summary}

The studies compose into a single argument about why formal methods stay niche and
what would change it. The barrier is the manual cost of specifications; inference
removes it, producing specifications that are accurate
(Chapter~\ref{ch:inference}) and valuable (Chapter~\ref{ch:testgen}); and inferred
specifications carried at the library boundary (Chapter~\ref{ch:distribution}),
because a caller's specification depends on its callees' (Chapter~\ref{ch:compositional}),
make verification across dependencies and behavioural-compatibility analysis ---
the latter a capability manual specification never made routine --- available to
ordinary software. The argument rests throughout on the disciplined separation of
heuristic inference from formal validation, the proposer-and-checker division that
lets the engine be aggressive and the verifier conservative. The principal
limitations are the homogeneity of the corpus, the use of a single model family,
the residual unverified clauses, and that the compatibility capability is
established as a mechanism rather than demonstrated end to end --- each named in the
chapter it most affects and each pointing toward the future work the next chapter
sets out.

The synthesis has aimed to show that the thesis is a single argument rather than
several. The argument is that the manual cost of specifications, long the binding
constraint on the adoption of formal methods, can be lifted by automatic inference
without sacrificing the soundness verification provides or the reach distribution
provides --- and that the resulting library-boundary specifications, propagating
from callee to caller, open behavioural-compatibility analysis as a new capability.
Each study removes one obstacle: inference removes the production cost, validation
removes the soundness objection, distribution removes the reach objection, and the
compositional measurement establishes the cross-library propagation the
compatibility capability depends on. The obstacles are removed not in the abstract
but by measurement on production code, and the residual limitations are the honest
record of where the removal is partial. What remains, after the studies, is an
argument that has survived the threats each was designed to expose it to, on the
evidence each was designed to gather --- which is the strongest form an empirical
thesis can take, and the form this one aspires to.
